% Generated draft: requires mathematical and source review.
\documentclass[11pt]{article}
\usepackage[margin=25mm]{geometry}
\usepackage{fontspec}
\setmainfont{Latin Modern Roman}
\usepackage{amsmath,amssymb,mathtools,unicode-math}
\setmathfont{Latin Modern Math}
\usepackage{enumitem,needspace,xurl,hyperref}
\setlength{\parindent}{0pt}
\setlength{\parskip}{5pt}
\setlength{\emergencystretch}{3em}
\newcommand{\N}{\mathbb{N}}
\newcommand{\Z}{\mathbb{Z}}
\newcommand{\Q}{\mathbb{Q}}
\newcommand{\R}{\mathbb{R}}
\newcommand{\C}{\mathbb{C}}
\newcommand{\F}{\mathbb{F}}
\newcommand{\A}{\mathbb{A}}
\newcommand{\PP}{\mathbb{P}}
\newcommand{\E}{\mathbb{E}}
\newcommand{\eps}{\varepsilon}
\newcommand{\dd}{\,\mathrm{d}}
\begin{document}
\section*{UnsolvedMath ID: 800103}
\textbf{Source problem code:} catalogue entry
\subsection{Definitions and mathematical statement}
For a foliation \(\mathcal F\) of the plane \(\R^2\), let
\[
L(\mathcal F)=\R^2/{\sim_{\mathcal F}},
\qquad x\sim_{\mathcal F}y\iff x\text{ and }y\text{ lie on the same leaf},
\]
with the quotient topology. The foliations in this question are understood to be nonsingular and one-dimensional. In the topological setting, two such foliations are transverse if they admit simultaneous local product charts in which their plaques are horizontal and vertical intervals. In a differentiable setting, transversality means that their tangent lines are distinct at every point. The supplied record does not specify the intended regularity category.

A possibly non-Hausdorff \(1\)-manifold is locally homeomorphic to \(\R\), without imposing the Hausdorff axiom. Here simply connected means path connected with trivial fundamental group. The local model \(\R\) uses the usual boundaryless interpretation; the source does not specify additional countability conventions.

\textbf{Pair realization problem.} Characterize the pairs \((\Lambda,\Lambda')\) of possibly non-Hausdorff simply connected \(1\)-manifolds for which there exist transverse foliations \(\mathcal F,\mathcal G\) of \(\R^2\) and homeomorphisms
\[
L(\mathcal F)\cong\Lambda,
\qquad L(\mathcal G)\cong\Lambda'.
\]
The regularity category and manifold conventions must be fixed to obtain a fully specified classification problem.

\textbf{Triple realization problem.} Characterize triples \((\Lambda_1,\Lambda_2,\Lambda_3)\) of such \(1\)-manifolds for which there exist foliations \(\mathcal F_1,\mathcal F_2,\mathcal F_3\) of \(\R^2\) with
\[
L(\mathcal F_i)\cong\Lambda_i\qquad(1\leq i\leq3),
\]
and with \(\mathcal F_i\) transverse to \(\mathcal F_j\) whenever \(i\ne j\). Pairwise transversality is the natural reading used here; the source only asks to replace ``pair'' by ``triple'' and does not explicitly define that condition.

\textbf{Anosov-flow specialization.} The source also asks for the analogous characterization of pairs of stable and unstable leaf spaces arising from Anosov flows. It further asks whether the stable leaf space determines the unstable leaf space. If determination means determination of the homeomorphism type, write \(L^s(\varphi)\) and \(L^u(\varphi)\) for these spaces for a flow \(\varphi\); the question would then be whether, for every two flows \(\varphi,\psi\) in the intended class,
\[
L^s(\varphi)\cong L^s(\psi)
\quad\Longrightarrow\quad
L^u(\varphi)\cong L^u(\psi).
\]
This last formulation is conditional: the supplied record does not specify the class of Anosov flows, whether the stable and unstable foliations are strong or weak, whether leaf spaces are taken on the original manifold, a universal cover, or an orbit space, or whether additional equivariant data are part of ``determines.'' These choices remain unresolved.
\subsection{Short English statement}
Determine which pairs, and which triples, of possibly non-Hausdorff simply connected \(1\)-manifolds occur as leaf spaces of transverse foliations of the plane. Also consider stable and unstable leaf spaces associated with Anosov flows, including whether the stable space determines the unstable space. Several conventions needed for a precise Anosov-flow formulation are absent from the supplied record. The supplied registry labels the problem open; this draft does not independently verify that status.
\subsection{Sources}
\begin{itemize}
\item Local UnsolvedMath registry, permanent ID \(800103\), source problem code \texttt{catalogue entry}, title ``Foliations and Diffeomorphism Groups --- Question 44.'' The supplied record belongs to the AMR Open Problem Lists collection and attributes the problem to Nariman and Yazdi, with proposed year \(2025\).
\item Supplied source citation: Nariman and Yazdi, \emph{Problem List on Foliations and Diffeomorphism Groups} (2025), Question 44, \url{https://arxiv.org/abs/2503.18274}. The supplied provenance records access on \texttt{2026-10-16} and extraction by \texttt{arxiv\_latex}; these are supplied provenance statements, not access or verification performed for this draft.
\item The supplied registry record, updated \texttt{2026-10-16}, labels the problem ``open.'' Its statement additionally reports that it is not yet known whether the stable leaf space determines the unstable leaf space. Both are supplied status assessments, not independently verified literature conclusions.
\end{itemize}
\end{document}
